% Minimal invoice - Pressa public template (MIT). https://pressa.dev/templates/invoice?design=minimal-invoice
% Line totals, subtotal, tax and total are computed from the data, each rounded to the
% currency's minor unit (cents, whole yen) as printed, so the printed numbers add up.
% Unit prices keep up to 4 decimals for items priced below a cent.
\documentclass[10pt,a4paper]{article}
\usepackage[T1]{fontenc}
\usepackage[scaled]{helvet}
\renewcommand{\familydefault}{\sfdefault}
\usepackage[margin=22mm,top=22mm]{geometry}
\usepackage[table]{xcolor}
\usepackage{tabularx}
\usepackage{longtable}
\usepackage{array}
\usepackage[hidelinks]{hyperref}

\definecolor{accent}{HTML}{0F766E}
\definecolor{ink}{HTML}{111111}
\definecolor{muted}{HTML}{777777}
\pagestyle{empty}
\setlength{\parindent}{0pt}
\newcolumntype{L}{>{\raggedright\arraybackslash}X}
\newcommand{\fieldname}[1]{\textcolor{muted}{\scriptsize\MakeUppercase{#1}}}

\begin{document}
\color{ink}
{% assign d = currency | currency_decimals %}{% assign subtotal = 0 %}{% for item in items %}{% assign price = item.unit_price | round: 4 %}{% assign line = item.quantity | times: price | round: d %}{% assign subtotal = subtotal | plus: line %}{% endfor %}{% assign tax = subtotal | times: tax_rate | divided_by: 100.0 | round: d %}{% assign total = subtotal | plus: tax %}
{\fontsize{30}{32}\selectfont Invoice}\par\vspace{8pt}\textcolor{accent}{\rule{12mm}{2pt}}

\vspace{22pt}
{\renewcommand{\arraystretch}{1.2}
\begin{tabularx}{\linewidth}{@{}L L L >{\raggedleft\arraybackslash}X@{}}
  \fieldname{Invoice no.} & \fieldname{Issued} & \fieldname{Due} & \fieldname{Amount due} \\
  {{ invoice_number }} & {{ issue_date }} & {{ due_date }} & \textcolor{accent}{\bfseries {{ total | currency: currency }}}
\end{tabularx}}

\vspace{22pt}
\begin{tabularx}{\linewidth}{@{}L@{\hspace{10mm}}L@{}}
  \fieldname{From} & \fieldname{Bill to} \\[2pt]
  {\bfseries\mbox{}{{ seller.name }}} & {\bfseries\mbox{}{{ client.name }}} \\
  {\small {{ seller.address }}\par {{ seller.email }}} & {\small {{ client.address }}}
\end{tabularx}

\vspace{26pt}
% A longtable, so a long list breaks across pages with the header repeated.
% Every column has a fixed width: one compile pass gives aligned columns.
{\renewcommand{\arraystretch}{1.7}\setlength{\LTpre}{0pt}\setlength{\LTpost}{0pt}\setlength{\LTleft}{0pt}\setlength{\LTright}{0pt}\arrayrulecolor{black!12}
\begin{longtable}{@{}>{\raggedright\arraybackslash}p{\dimexpr\linewidth-6\tabcolsep-12mm-26mm-26mm\relax} >{\raggedleft\arraybackslash}p{12mm} >{\raggedleft\arraybackslash}p{26mm} >{\raggedleft\arraybackslash}p{26mm}@{}}
  \fieldname{Description} & \fieldname{Qty} & \fieldname{Unit price} & \fieldname{Amount} \\
  \hline
\endhead
{% for item in items %}{% assign price = item.unit_price | round: 4 %}{% assign line = item.quantity | times: price | round: d %}  {{ item.description }} & {{ item.quantity | round: 2 }} & {{ price | currency: currency, 4 }} & {{ line | currency: currency }} \\
  \hline
{% endfor %}\end{longtable}}

\vspace{10pt}
\hfill{\renewcommand{\arraystretch}{1.5}\begin{tabular}{@{}l@{\hspace{14mm}}r@{}}
  \textcolor{muted}{Subtotal} & {{ subtotal | currency: currency }} \\
  \textcolor{muted}{Tax ({{ tax_rate }}\%)} & {{ tax | currency: currency }} \\
  \textbf{Total} & \textcolor{accent}{\large\bfseries {{ total | currency: currency }}} \\
\end{tabular}}

\vfill
\begin{tabularx}{\linewidth}{@{}L@{\hspace{10mm}}L@{\hspace{10mm}}L@{}}
  \fieldname{Terms} & \fieldname{Payment} & \fieldname{Notes} \\[2pt]
  {\small {{ payment_terms }}} & {\small {{ payment_details }}} & {\small {{ notes }}}
\end{tabularx}
\end{document}
